% =============================================================================
% Chapter 4 — Controller Design and Calibration
% Tara Torbati — Master's Thesis
%
% This chapter introduces the four controller architectures evaluated in
% Chapter 5 and documents all design and calibration decisions required to
% reproduce the experimental results. The MPC formulation, the five-term
% cost function, the symbolic NLP construction, and the four-phase
% weight-sensitivity analysis are developed in detail. The SAC architecture
% is then presented with its corrected MDP formulation, CTDE actor-critic
% structure, training protocol, and the AR(1) forecast-noise model shared
% with the MPC robustness evaluation.
% =============================================================================

\chapter{Controller Design and Calibration}
\label{chap:controllers}

This chapter introduces the four controller architectures evaluated on the
shared virtual plant of Chapter~\ref{chap:system}: a no-irrigation rainfed
baseline, a fixed-schedule heuristic representing traditional Gilan practice,
a constrained Model Predictive Controller (MPC), and a Soft Actor-Critic
(SAC) reinforcement learning agent. The two baselines are introduced briefly
in Section~\ref{sec:ctrl-baselines} and serve to bracket the performance
space within which the optimization-based controllers operate.

The MPC is developed in detail across
Sections~\ref{sec:ctrl-mpc-formulation}--\ref{sec:ctrl-recommended}. This
development spans the multi-objective cost function with economic anchoring
(Section~\ref{sec:ctrl-mpc-formulation}), the Control-Oriented Model and
symbolic NLP construction that make online optimization computationally
tractable (Section~\ref{sec:ctrl-com}), the prediction horizon selection
(Section~\ref{ssec:ctrl-horizon}), the four-phase weight-sensitivity analysis
that identifies and resolves three cost-function pathologies
(Section~\ref{sec:ctrl-sweep}), and the resulting recommended operating
point (Section~\ref{sec:ctrl-recommended}).

The SAC architecture is presented in Section~\ref{sec:ctrl-sac}, covering
the Markov Decision Process formulation, the Gymnasium environment wrapper,
the CTDE actor-critic structure, and the full training protocol including
the 20/3/3 train/dev/test year split. Section~\ref{sec:ctrl-forecast-noise}
introduces the AR(1) meteorological forecast noise model used to evaluate
both controllers under realistic forecast uncertainty in Chapter~\ref{chap:results}.

% =============================================================================
\section{Baselines}
\label{sec:ctrl-baselines}
% =============================================================================

\subsection{No-Irrigation Baseline}
\label{ssec:ctrl-noirr}

The no-irrigation controller applies $u^{(n)}(k) = 0$ for all agents and all
timesteps. The crop subsists on rainfall alone. This baseline establishes the
rainfed yield floor against which any irrigation controller is evaluated.
Under the dry scenario (2022, 39.7~mm seasonal rainfall), the no-irrigation
controller produces a near-complete yield collapse of 1{,}462~kg/ha with
87~drought-days per agent. Under the wet scenario (2024, 176.8~mm), partial
yield of 2{,}243~kg/ha is maintained with 83~drought-days per agent. These
values quantify the marginal value of irrigation across scenarios and provide
the lower bound for the comparative analysis of Chapter~\ref{chap:results}.

\subsection{Fixed-Schedule Heuristic}
\label{ssec:ctrl-fixed}

The fixed-schedule controller represents traditional Hashemi irrigation
practice in Gilan~\cite{jalali2021}. The seasonal water budget
$W_{\mathrm{total}}$ is divided uniformly across all $K = 93$ days and
applied identically to all $N = 130$ agents:
\begin{equation}
u_{\mathrm{fixed}}^{(n)}(k) = \frac{W_{\mathrm{total}}}{K}
\qquad \forall n,\, \forall k
\label{eq:fixed-schedule}
\end{equation}
This open-loop policy uses the full budget exactly, distributes it evenly
in space and time, and ignores both spatial heterogeneity and temporal
climate variation. As a baseline it isolates the value of \emph{closed-loop
feedback}: any improvement of MPC or SAC over the fixed schedule is
attributable specifically to the controller's ability to react to state and
forecast information, since all three controllers share the same ABM, the
same budget, and the same actuator bounds.

Under the dry scenario at full budget (484~mm), the fixed schedule yields
3{,}607~kg/ha but incurs 54~waterlog-days and 28~drought-days per agent: the
uniform application is simultaneously too much (waterlogging upslope low-ET
agents) and too little (under-supplying high-ET periods). Under the wet
scenario the waterlogging worsens further (83.7~waterlog-days) because the
fixed schedule does not react to rainfall, delivering the same volume
regardless of whether the crop already has sufficient moisture from
precipitation.

% =============================================================================
\section{MPC Problem Formulation}
\label{sec:ctrl-mpc-formulation}
% =============================================================================

The Model Predictive Controller solves the constrained finite-horizon optimal
control problem of Section~\ref{sec:sys-control} in a receding-horizon
fashion. At each timestep~$k$, the controller observes the field state
$\mathbf{X}(k)$, retrieves the climate forecast over the next $H_p$ days,
solves a nonlinear program (NLP) to obtain the optimal control sequence
$\mathbf{U}^*(k:k+H_p-1)$, applies only the first-step action
$\mathbf{u}^*(k)$, and re-solves at $k+1$ with updated state and forecast.
This receding-horizon mechanism converts the open-loop optimization into a
feedback control law and provides robustness against modeling errors and
weather-forecast deviations.

The per-iteration optimization is
\begin{equation}
\min_{\mathbf{U}(k:k+H_p-1)}
J\bigl(\mathbf{X}(k:k+H_p),\,\mathbf{U}(k:k+H_p-1)\bigr)
\label{eq:mpc-ocp}
\end{equation}
subject to
\begin{align}
\mathbf{x}^{(n)}(j+1) &= f\!\bigl(\mathbf{x}^{(n)}(j),\,u^{(n)}(j),\,
  \mathbf{w}(j);\,\mathcal{G}\bigr),
  \quad n=1,\ldots,N,\quad j=k,\ldots,k+H_p-1
  \label{eq:mpc-dyn}\\
0 &\leq u^{(n)}(j) \leq u_{\max}, \qquad \forall n,\,\forall j
  \label{eq:mpc-bounds}\\
\sum_{j=k}^{k+H_p-1}\sum_{n=1}^{N} u^{(n)}(j) &\leq W_{\mathrm{remaining}}(k)
  \label{eq:mpc-budget}
\end{align}
where $W_{\mathrm{remaining}}(k)$ is the water budget remaining at the start
of timestep~$k$. The remainder of this section addresses three components of
this formulation that require careful design: the cost function $J$
(Section~\ref{ssec:ctrl-cost}), the symbolic Control-Oriented Model and
NLP construction (Section~\ref{sec:ctrl-com}), and the choice of prediction
horizon $H_p$ (Section~\ref{ssec:ctrl-horizon}).

\subsection{Multi-Objective Cost Function}
\label{ssec:ctrl-cost}

The cost function is the principal design lever of the MPC. It encodes the
agronomic and economic priorities traded off when allocating water across
space and time. The formulation adopted here is a five-term sum:
\begin{equation}
J =
  \underbrace{-\alpha_1 \cdot J_{\mathrm{biomass}}}_{\text{terminal reward}}
+ \underbrace{\alpha_2 \cdot J_{\mathrm{water}}}_{\text{economic cost}}
+ \underbrace{\alpha_3 \cdot J_{\mathrm{drought}}}_{\text{drought penalty}}
+ \underbrace{\alpha_5 \cdot J_{\Delta u}}_{\text{actuator regularizer}}
+ \underbrace{\alpha_6 \cdot J_{\mathrm{overFC}}}_{\text{FC-overshoot penalty}}
\label{eq:cost-fn}
\end{equation}
A sixth term $\alpha_4 J_{\mathrm{ponding}}$ appeared in an initial
formulation and is set to zero at the recommended operating point; its
empirical falsification is the subject of
Sections~\ref{ssec:ctrl-cost-overfc} and \ref{ssec:ctrl-sweep-joint}.

\subsubsection*{Term 1 — Terminal Biomass}

The terminal biomass reward is the principal positive driver of the cost:
\begin{equation}
J_{\mathrm{biomass}} = \frac{1}{N \cdot x_{4,\mathrm{ref}}}
  \sum_{n=1}^{N} x_4^{(n)}(K)
\label{eq:J-biomass}
\end{equation}
where $x_{4,\mathrm{ref}} = 1{,}000$~g/m$^2$ normalizes to unit order of
magnitude at full agronomic potential. Within the receding-horizon loop, the
end-of-horizon state $\mathbf{X}(k+H_p)$ replaces $\mathbf{X}(K)$,
following standard Mayer-term MPC convention.

\subsubsection*{Term 2 — Water Consumption (Economic Anchoring)}

The water-cost term penalizes cumulative volumetric water applied over the
horizon:
\begin{equation}
J_{\mathrm{water}} = \frac{1}{N \cdot u_{\max} \cdot H_p}
  \sum_{j=k}^{k+H_p-1} \sum_{n=1}^{N} u^{(n)}(j)
\label{eq:J-water}
\end{equation}
Unlike the other weights, $\alpha_2$ admits direct economic interpretation.
Iranian water policy is structured around a four-tier volumetric tariff
schedule~\cite{Mesgaran2018}. Setting $\alpha_2$ proportional to a chosen
tariff tier allows the controller's behaviour to be interpreted as a response
to a real economic price signal rather than an abstract regularization.
The four operationally relevant tiers and their corresponding $\alpha_2$
values are:
\begin{itemize}
  \item Subsidized agricultural ($\sim$175~Toman/m$^3$):
    $\alpha_2 = 0.0004$
  \item Domestic base ($\sim$7{,}000~Toman/m$^3$):
    $\alpha_2 = 0.016$
  \item Domestic high-consumption ($\sim$19{,}000~Toman/m$^3$):
    $\alpha_2 = 0.044$
  \item Industrial ($\sim$115{,}000~Toman/m$^3$):
    $\alpha_2 = 0.265$
\end{itemize}
The recommended value $\alpha_2^* = 0.016$ corresponds to the domestic-base
tier; the empirical justification is presented in
Section~\ref{ssec:ctrl-sweep-a2}.

\subsubsection*{Term 3 — Drought Stress}

The drought-stress term penalizes per-day soil-moisture deficits below the
agronomic stress threshold $\mathrm{ST} = \mathrm{FC}(1 - p)$, where $p = 0.20$
is the FAO depletion fraction for rice~\cite{allen1998}:
\begin{equation}
J_{\mathrm{drought}} = \frac{1}{N \cdot H_p \cdot (\mathrm{ST} - \mathrm{WP})}
  \sum_{j=k}^{k+H_p-1} \sum_{n=1}^{N}
  \mathrm{smax}_{\varepsilon}\!\bigl(\mathrm{ST} - x_1^{(n)}(j),\, 0\bigr)
\label{eq:J-drought}
\end{equation}
The smooth-max operator $\mathrm{smax}_{\varepsilon}(a,b) =
\tfrac{1}{2}\bigl(a + b + \sqrt{(a-b)^2 + \varepsilon^2}\bigr)$ with
$\varepsilon = 10^{-2}$ replaces the non-differentiable $\max(\cdot, 0)$,
enabling second-order gradient information for the IPOPT solver while
preserving $C^2$ continuity everywhere on $\mathbb{R}$. The denominator
normalizes the term to $[0,1]$ when the soil is held continuously at the
wilting point.

\subsubsection*{Term 5 — Actuator Smoothing}

Sudden swings in irrigation flow stress hardware and complicate operational
scheduling. A quadratic penalty on the daily increment
$\Delta u^{(n)}(j) = u^{(n)}(j) - u^{(n)}(j-1)$ smooths the control
trajectory:
\begin{equation}
J_{\Delta u} = \frac{1}{N \cdot u_{\max}^2 \cdot H_p}
  \sum_{j=k}^{k+H_p-1} \sum_{n=1}^{N}
  \bigl(\Delta u^{(n)}(j)\bigr)^2
\label{eq:J-deltau}
\end{equation}
The recommended value $\alpha_5^* = 0.005$ is small relative to the other
path penalties; the term acts as a regularizer that breaks ties among
otherwise equivalent control sequences without materially shifting the
optimum.

\subsubsection*{Term 6 — Field-Capacity Overshoot Penalty}
\label{ssec:ctrl-cost-overfc}

The two soft state penalties in \eqref{eq:cost-fn} are $J_{\mathrm{drought}}$
(soil moisture below the stress threshold) and $J_{\mathrm{overFC}}$
(soil moisture above field capacity). The latter addresses a controller
pathology that emerged empirically from nominal MPC runs. The mechanism is
described here; the supporting evidence is presented in
Sections~\ref{ssec:ctrl-sweep-a4} and~\ref{ssec:ctrl-sweep-joint}.

The initial cost-function formulation included an explicit ponding penalty on
the surface state $x_5$:
\begin{equation*}
J_{\mathrm{ponding}} = \frac{1}{N \cdot x_{5,\mathrm{ref}} \cdot H_p}
  \sum_{j=k}^{k+H_p-1} \sum_{n=1}^{N} x_5^{(n)}(j)
\end{equation*}
The motivation was that surface ponding causes root-zone oxygen deprivation,
and $x_5$ and $x_1$ describe physically distinct hydrological compartments.
Empirically, this expectation was falsified. Nominal runs at $H_p = 14$
revealed that the controller systematically pushed water past the surface
(suppressing $x_5$) into the root zone, where $x_1$ persistently exceeded
field capacity for 21.6~agent-days at dry/100\%/$H_p = 14$ and
61.1~agent-days at wet/100\%/$H_p = 14$, with peak $x_1$ reaching 205~mm.
The agronomic consequence of $x_1 > \mathrm{FC}$ is root hypoxia within
24--48~hours~\cite{setter1997}; the ponding penalty did not prevent this
because it was targeting the wrong state variable.

The diagnosis motivated a sixth cost term targeting $x_1 > \mathrm{FC}$
directly:
\begin{equation}
J_{\mathrm{overFC}} = \frac{1}{N \cdot \mathrm{FC}^2 \cdot H_p}
  \sum_{j=k}^{k+H_p-1} \sum_{n=1}^{N}
  \mathrm{smax}_{\varepsilon}\!\bigl(x_1^{(n)}(j) - \mathrm{FC},\, 0\bigr)^2
\label{eq:J-overfc}
\end{equation}
The quadratic form makes the penalty $C^2$-smooth and ensures that large
overshoots are penalized strongly while small transient excursions are
penalized lightly. Whether the original $J_{\mathrm{ponding}}$ term retains
any independent agronomic role once $\alpha_6 > 0$ is tested in the joint
factorial of Section~\ref{ssec:ctrl-sweep-joint}; the conclusion is that at
sufficient $\alpha_6$ strength the ponding penalty contributes no measurable
benefit and the cost function reduces to the five-term form
\eqref{eq:cost-fn}. The $J_{\mathrm{ponding}}$ term is retained in the
implementation but inactive at $\alpha_4 = 0$, leaving the formulation
extensible.

% =============================================================================
\section{Control-Oriented Model and NLP Construction}
\label{sec:ctrl-com}
% =============================================================================

The optimization \eqref{eq:mpc-ocp}--\eqref{eq:mpc-budget} is implemented
in CasADi~\cite{andersson2019casadi}, a symbolic automatic-differentiation
framework that constructs a single NLP for the entire prediction horizon.
Three design decisions reduce the problem to a tractable size and enable
efficient gradient evaluation.

\subsection{State-Vector Reduction}
\label{ssec:ctrl-state-reduction}

Of the five ABM state variables, only $x_1$ (root-zone moisture) and $x_5$
(surface ponding) are directly affected by the control input within a single
timestep. The remaining states evolve from climate forcings alone:
\begin{itemize}
  \item $x_2$ (cumulative thermal time) is a function of temperature only,
    deterministic over any forecast window.
  \item $x_3$ (maturation index) accumulates $\theta_{11}(1-h_2)
    + \theta_{12}(1-h_3)$ per day; both $h_2$ (heat stress) and $h_3$
    (drought stress) are scalar functions of temperature and $x_1$,
    respectively.
  \item $x_4$ (biomass) accumulates as a function of $x_1$, $x_2$, $x_3$,
    and shortwave radiation.
\end{itemize}
Within the NLP, $x_2$, $x_3$, and $x_4$ are computed by forward integration
as symbolic expressions of the shooting states $x_1$ and $x_5$ rather than
carried as independent decision variables. This reduces the shooting-state
dimension from $5 \cdot N \cdot H_p$ to $2 \cdot N \cdot H_p$ at the cost
of slightly more complex symbolic expressions. At $N = 130$ and $H_p = 8$,
the NLP has $2 \cdot 130 \cdot 8 = 2{,}080$ state variables plus
$130 \cdot 8 = 1{,}040$ action variables, for a total of \textbf{3,120
decision variables} and \textbf{2,081 constraints} (1,040 bound constraints
plus the single budget constraint at each of the 8 horizon steps, plus the
2,080 state-transition equalities, minus double counting of shared horizon
endpoints).

\subsection{Deterministic Caching of Biological Nonlinearities}
\label{ssec:ctrl-caching}

The growth function $g(x_2)$, the heat-stress factor $h_2$, the
cold-stress factor $h_7$, and the crop atmospheric demand $K_c \cdot
\mathrm{ET}_0$ (denoted $\mathrm{ET}_c$) depend only on temperature,
radiation, and thermal time --- all of which are deterministic functions of
the forecast rather than of the irrigation action. These quantities are
therefore precomputed over the full season for each scenario (or, in the
randomized training mode, for each sampled historical year) and cached on
disk as NumPy arrays. Within the NLP, they enter as fixed numeric constants
rather than as symbolic expressions, substantially reducing the symbolic
graph complexity and accelerating both function evaluations and gradient
computations.

Specifically, for each day $d$ of the season, the following quantities are
precomputed:
\begin{itemize}
  \item $h_1^{(d)} = \max(T_{\mathrm{mean}}^{(d)} - \theta_7,\, 0)$,
    the daily thermal increment;
  \item $x_2^{(d)} = x_2(0) + \sum_{i=0}^{d-1} h_1^{(i)}$,
    the cumulative thermal time;
  \item $h_2^{(d)}$, the heat-stress factor;
  \item $h_7^{(d)}$, the low-temperature stress factor;
  \item $g_{\mathrm{base}}^{(d)} = g(x_2^{(d)})$ evaluated at
    $x_3 = 0$ (the no-maturation-stress baseline);
  \item $\mathrm{ET}_c^{(d)} = K_c \cdot \mathrm{ET}_0^{(d)}$,
    the crop-specific atmospheric demand.
\end{itemize}
These six arrays are used both by the MPC (as NLP constants) and by the SAC
observation builder (as features in the 707-dimensional observation vector),
ensuring that both controllers operate on identical agronomic context signals.

\subsection{Symbolic Topological Cascade Unrolling}
\label{ssec:ctrl-cascade-unrolling}

The runoff-routing graph $\mathcal{G}$ is statically known at NLP
construction time. The agents are sorted topologically (highest elevation
first) and the cascade is unrolled symbolically: each agent's surface-water
balance is expressed as an explicit function of its own action $u^{(n)}(j)$
and of the runoff contributions of its uphill neighbors, which have already
been computed in the topological sweep. This produces a single forward pass
per timestep with no implicit equations or iterative routing solvers. The
resulting NLP is dense in the spatial dimension (each agent's state depends
on all uphill agents' actions) but the symbolic graph is straight-line code
that CasADi's automatic differentiation processes efficiently.

\subsection{Smooth Approximations for Solver Convergence}
\label{ssec:ctrl-smooth}

The ABM contains several non-differentiable operations that must be smoothed
for interior-point optimization. The smooth-max operator introduced in
Section~\ref{ssec:ctrl-cost} ($\mathrm{smax}_\varepsilon$ with
$\varepsilon = 0.01$) is applied uniformly wherever $\max(\cdot, 0)$,
$\min(\cdot, \cdot)$, or clamp-to-bounds operations appear in the cost
function and state-transition equations. The $C^2$-smooth approximation
preserves the function value to within $\varepsilon/2 = 0.005$~mm across the
relevant operating range and provides the twice-differentiable structure
required by IPOPT's filter line-search globalization strategy.

\subsection{NLP Solver Configuration}
\label{ssec:ctrl-solver}

The complete NLP is solved with IPOPT~\cite{wachter2006ipopt} at each
receding-horizon step. The principal solver settings are:
\begin{itemize}
  \item Convergence tolerance: $\mathrm{tol} = 10^{-4}$
  \item Acceptable tolerance: $\mathrm{acceptable\_tol} = 10^{-3}$
  \item Maximum iterations: $500$
  \item Warm-starting: the action sequence from the previous solve, shifted
    by one timestep (last entry zero-initialized), is used as the
    initial guess for the current solve (Section~\ref{ssec:ctrl-warmstart}).
\end{itemize}
These settings balance solver reliability against per-iteration latency. The
relaxed primary tolerance is justified by the observation that the cost
surface near the field-capacity saturation regime has small agronomic
gradients, so a tight $10^{-8}$ tolerance produces only marginal changes
in the optimal action while requiring substantially longer solve times.

\paragraph{Observed convergence.}
At the recommended operating point $\boldsymbol{\alpha}^*$ and $H_p = 8$,
the mean per-day solve time across the six evaluated scenario/budget
combinations is 25.9~s (median 22.1~s), with a worst-case of 274~s occurring
in a high-rainfall wet-year iteration where the constraint surface is
particularly complex. The total wall-clock time per 93-day season ranges
from 22 to 51~minutes depending on scenario and budget, confirming practical
feasibility for daily-dispatch agricultural systems. At $H_p = 14$, mean
solve time rises to 292~s (5-fold increase) with a worst-case of 12{,}134~s
in the wet/85\% configuration, a latency that exceeds practical dispatch
windows for any system faster than daily scheduling.

\subsection{Warm-Starting and Computational Acceleration}
\label{ssec:ctrl-warmstart}

At each timestep $k$, the initial guess for the NLP is constructed by
taking the optimal action sequence from step $k-1$, removing the first
element (which was applied), and appending a zero-action as the final entry.
This shifted warm-start provides a feasible (or near-feasible) initial
point that lies close to the true optimum by continuity of the optimal
value function, reducing the number of IPOPT iterations required per solve.
In practice, warm-started solves at $H_p = 8$ converge in a mean of
112~iterations versus approximately 280~iterations from a cold start, a
$2.5\times$ reduction in per-step compute.

\subsection{Prediction Horizon}
\label{ssec:ctrl-horizon}

Three prediction horizons are evaluated: $H_p \in \{3, 8, 14\}$~days. The
trade-off is fundamental. Short horizons produce small NLPs that solve in
under one second ($H_p = 3$: mean 0.8~s) and respond quickly to forecast
updates, but cannot anticipate multi-day rainfall events or stage water
deliveries against a multi-week phenological schedule. Long horizons enable
more strategic allocation but produce NLPs whose worst-case solve times
exceed three hours at $H_p = 14$ in wet scenarios, and whose larger search
space increases susceptibility to non-convex local-minimum traps.

The choice of $H_p$ is informed by the empirical analysis of
Section~\ref{sec:ctrl-recommended}. To anticipate that result briefly: under
the original cost function ($\alpha_6 = 0$), $H_p = 14$ exhibited a
horizon-inversion pathology in which yields were strictly worse than at
shorter horizons, traceable to the field-capacity overshoot problem of
Section~\ref{ssec:ctrl-cost-overfc}. Under the recommended cost function,
this pathology is resolved and yields converge across horizons. The
recommended horizon is $H_p^* = 8$, selected on the basis that yields are
statistically equivalent across horizons at $\boldsymbol{\alpha}^*$ while
per-step solve time at $H_p = 8$ is approximately $5\times$ faster than at
$H_p = 14$ and the warm-start mechanism is most effective in the mid-range
horizon.

% =============================================================================
\section{Weight-Sensitivity Analysis}
\label{sec:ctrl-sweep}
% =============================================================================

This section presents the systematic sensitivity sweep that calibrated the
cost-function weights $\boldsymbol{\alpha} = (\alpha_1, \ldots, \alpha_6)$
and identified the recommended operating point. The sweep was conducted in
four phases across 33 controller configurations and approximately 70~hours
of compute. The four phases are:
\begin{enumerate}
  \item \textbf{One-at-a-time (OAT) sweep} of the four weights with
    ambiguous economic or agronomic motivation ($\alpha_2, \alpha_3,
    \alpha_4, \alpha_6$), Sections~\ref{ssec:ctrl-sweep-a2}--\ref{ssec:ctrl-sweep-a6}.
  \item \textbf{Joint $(\alpha_4, \alpha_6)$ factorial} at the worst-case
    wet/$H_p=14$ stress-test point, testing the redundancy hypothesis between
    the two waterlogging penalties, Section~\ref{ssec:ctrl-sweep-joint}.
  \item \textbf{Cross-scenario validation} of two candidate operating points
    across the operating-point grid, Section~\ref{ssec:ctrl-sweep-validation}.
  \item \textbf{Ceiling test} confirming that the recommended $\alpha_6 = 8$
    sits at the asymptotic plateau of the weight-vs-performance response,
    Section~\ref{ssec:ctrl-sweep-ceiling}.
\end{enumerate}

The nominal reference against which sweeps are compared is
$\boldsymbol{\alpha}^{\mathrm{nom}} = (1.0, 0.01, 0.1, 0.5, 0.005, 0)$.
This reference deactivates $J_{\mathrm{overFC}}$ entirely ($\alpha_6 = 0$)
and uses an intermediate water-cost weight ($\alpha_2 = 0.01$), so that the
price-tier sweep is not biased toward any specific tariff.

\paragraph{Diagnostic metrics.} Each configuration is evaluated on the
following metrics: terminal yield (kg/ha), total water applied (mm),
waterlog-days per agent (days with $x_1 > \mathrm{FC} = 140$~mm),
severe-waterlog-days per agent (days with $x_1 > 180$~mm, the root-hypoxia
threshold), peak $x_1$ (mm), and worst-case per-day solve time (s).

\subsection{\texorpdfstring{$\alpha_2$: Water-Price Tiers}{Alpha2: Water-Price Tiers}}
\label{ssec:ctrl-sweep-a2}

The $\alpha_2$ sweep evaluates the controller's response to the four Iranian
water-pricing tiers at dry/100\%/$H_p = 8$. Results are summarized in
Table~\ref{tab:sweep-a2}.

\begin{table}[h]
\centering
\caption{$\alpha_2$ sweep: dry/100\%/$H_p=8$. Yield, water use, and
waterlog-days as functions of the water-cost weight, mapped to Iranian
water-pricing tiers.}
\label{tab:sweep-a2}
\begin{tabular}{lcccc}
\toprule
\textbf{Tier} & $\boldsymbol{\alpha_2}$
  & \textbf{Yield (kg/ha)} & \textbf{Water (mm)} & \textbf{wlog-d} \\
\midrule
Subsidized agricultural & 0.0004 & 3{,}836 & 482.4 & 76.0 \\
Domestic base           & 0.016  & 4{,}142 & 468.2 &  1.8 \\
Domestic high           & 0.044  & 4{,}131 & 465.6 &  0.0 \\
Industrial              & 0.265  & 1{,}495 &   6.0 &  0.0 \\
\bottomrule
\end{tabular}
\end{table}

Three findings emerge. First, at the subsidized agricultural rate, the
controller over-irrigates: water use is the highest in the sweep (482.4~mm)
yet yield is the lowest (3{,}836~kg/ha) due to severe field-capacity
overshoot — 76~waterlog-days, an order of magnitude worse than any other
configuration. Cheap water actively damages the crop because the controller
has no economic incentive to refrain from over-application. Second, at the
domestic-base and domestic-high tiers, the controller curbs water use
marginally while achieving yields above 4{,}100~kg/ha with negligible
waterlogging. The moderate price discipline forces the controller into the
agronomic sweet spot. Third, at the industrial tier, the controller adopts a
rational-shutdown policy: water use collapses to 6~mm and yield drops to
1{,}495~kg/ha, because the marginal water cost exceeds the marginal biomass
revenue at the industrial price level. This response is economically correct:
at industrial pricing, rice cultivation is no longer profitable.

The recommended value $\alpha_2^* = 0.016$ corresponds to the domestic-base
tier, which achieves the best yield-per-water trade-off without inducing
either the over-application pathology or the rational-shutdown collapse.

\subsection{\texorpdfstring{$\alpha_3$: Drought Stress Regularizer}{Alpha3: Drought Stress}}
\label{ssec:ctrl-sweep-a3}

The $\alpha_3$ sweep evaluates values of $\{0.03, 0.30\}$ at
dry/100\%/$H_p = 8$, spanning a full order of magnitude around the nominal
$\alpha_3 = 0.1$. Yield variations are below 1\% and water-use variations
are below 1\%. The drought-stress term is agronomically redundant under the
tested conditions because the implicit drought penalty embedded in the
biomass increment (the $h_3$ multiplier in $\dot{x}_4$) already provides a
strong gradient against drought. The explicit $\alpha_3$ term retains
diagnostic utility under tight budget conditions: at 70\% budget where the
controller cannot avoid drought, higher $\alpha_3$ values produce strongly
positive ex-post cost values that flag the infeasibility to an operator.
The default $\alpha_3 = 0.1$ is retained.

\subsection{\texorpdfstring{$\alpha_4$: Surface-Ponding Penalty (Falsification)}{Alpha4: Surface Ponding}}
\label{ssec:ctrl-sweep-a4}

The $\alpha_4$ sweep evaluates values of $\{2.0, 5.0\}$ (versus nominal
$0.5$) at wet/100\%/$H_p=14$, the operating regime where surface ponding is
most prevalent. Results are summarized in Table~\ref{tab:sweep-a4}.

\begin{table}[h]
\centering
\caption{$\alpha_4$ sweep at wet/100\%/$H_p=14$ with $\alpha_6 = 0$.
Waterlog-days increase before decreasing; root-zone overshoot is
unaffected: $\alpha_4$ alone does not resolve the wet-scenario pathology.}
\label{tab:sweep-a4}
\begin{tabular}{ccccc}
\toprule
$\boldsymbol{\alpha_4}$ & \textbf{Yield}
  & \textbf{Water (mm)} & \textbf{wlog-d} & $x_5^{\mathrm{max}}$~\textbf{(mm)} \\
\midrule
0.5 (nominal) & 3{,}392 & 446.0 & 61.1 & 40.8 \\
2.0           & 3{,}334 & 475.1 & 68.1 &  0.7 \\
5.0           & 3{,}537 & 407.8 & 47.0 & 0.04 \\
\bottomrule
\end{tabular}
\end{table}

At $\alpha_4 = 2.0$, surface ponding is essentially eliminated
($x_5^{\mathrm{max}}$ collapses from 40.8 to 0.7~mm), but root-zone
waterlog-days \emph{increase} from 61 to 68. The controller has displaced
the pathology from the surface to the subsurface: water is pushed past the
surface and into the root zone where $\alpha_4$ cannot see it. Only at
$\alpha_4 = 5.0$ does behaviour improve, and only by suppressing total water
application (water use drops 38~mm) rather than by improved scheduling.
The conclusion is that $\alpha_4$ alone cannot resolve the wet-scenario
pathology because it targets $x_5$ instead of $x_1$.

\subsection{\texorpdfstring{$\alpha_6$: Field-Capacity Overshoot}{Alpha6: FC Overshoot}}
\label{ssec:ctrl-sweep-a6}

The $\alpha_6$ sweep evaluates three values ($\{0.1, 0.5, 2.0\}$) at
dry/$H_p = 8$ and one value ($\alpha_6 = 2.0$) at the horizon-inversion
regime dry/$H_p = 14$. At $H_p = 8$, the FC-overshoot is small at the
nominal weights (1.9~waterlog-days), and sweep values produce yield
variations within 1\%. The penalty has nothing to act against in this regime,
confirming that $\alpha_6$ is a horizon-conditional lever rather than a
universally beneficial one.

At $H_p = 14$, $\alpha_6 = 2.0$ reduces dry/100\% waterlog-days from 21.6
to 13.9 (35\% reduction) and improves yield slightly. The mechanism works as
designed at the horizon where the pathology is active. However, 13.9
residual waterlog-days remain approximately $7\times$ the $H_p = 8$ baseline,
indicating that $\alpha_6 = 2.0$ is insufficient, which motivates the broader
joint sweep with higher $\alpha_6$ values.

\subsection{\texorpdfstring{Joint $(\alpha_4, \alpha_6)$ Factorial}{Joint Alpha4 Alpha6 Factorial}}
\label{ssec:ctrl-sweep-joint}

The joint factorial tests whether $\alpha_4$ retains independent agronomic
value once $\alpha_6$ is active. The factorial spans $\alpha_4 \in \{0, 0.5,
2, 5\}$ and $\alpha_6 \in \{0, 2, 4, 8\}$ across nine new configurations at
wet/100\%/$H_p=14$, combined with three preexisting runs at $\alpha_6 = 0$.
Results are summarized in Table~\ref{tab:sweep-joint}.

\begin{table}[h]
\centering
\caption{Joint $(\alpha_4, \alpha_6)$ factorial at wet/100\%/$H_p=14$.
Sev-wlog reports days with $x_1 > 180$~mm. Sorted by waterlog-day count.}
\label{tab:sweep-joint}
\begin{tabular}{cccccc}
\toprule
$\boldsymbol{\alpha_4}$ & $\boldsymbol{\alpha_6}$
  & \textbf{Yield} & \textbf{Water} & \textbf{wlog-d} & \textbf{sev-wlog} \\
\midrule
0   & 8 & 3{,}774 & 321.5 & 21.1 & 0.0 \\
5   & 4 & 3{,}756 & 319.5 & 26.3 & 0.0 \\
2   & 4 & 3{,}763 & 325.9 & 28.2 & 0.0 \\
0   & 4 & 3{,}677 & 357.8 & 29.8 & 4.8 \\
0.5 & 4 & 3{,}749 & 330.5 & 31.4 & 1.0 \\
5   & 2 & 3{,}680 & 343.0 & 35.0 & 3.3 \\
0   & 2 & 3{,}684 & 361.8 & 36.0 & 4.8 \\
2   & 2 & 3{,}659 & 347.4 & 41.8 & 5.1 \\
0   & 0 & 3{,}294 & 478.4 & 63.7 & 24.3 \\
\bottomrule
\end{tabular}
\end{table}

Three findings determine the recommended cost function. First, the
configuration $(\alpha_4 = 0, \alpha_6 = 8)$ achieves the highest yield,
lowest waterlog-days, and zero severe waterlogging without any explicit
ponding penalty. The mechanism is dynamical: surface water must eventually
either infiltrate (raising $x_1$, penalized by $\alpha_6$) or drain
off-farm. A sufficiently strong $\alpha_6$ therefore indirectly suppresses
high-$x_5$ states by making the downstream infiltrated state expensive.
Second, at intermediate $\alpha_6 = 4$, configurations with $\alpha_4 \geq 2$
eliminate severe waterlogging while $\alpha_4 = 0$ leaves 4.8~days, showing
that $\alpha_4$ retains marginal utility in the intermediate regime.
Third, the $(\alpha_4 = 0, \alpha_6 = 0)$ baseline differs from the nominal
$(\alpha_4 = 0.5, \alpha_6 = 0)$ by less than solver noise, demonstrating
that the original $\alpha_4 = 0.5$ setting was operationally inactive.

The factorial identifies two candidate operating points for cross-scenario
validation: \textbf{Option~A} $(\alpha_4 = 0, \alpha_6 = 8)$, the
parsimonious five-term cost function, and \textbf{Option~B}
$(\alpha_4 = 2, \alpha_6 = 4)$, the joint six-term cost function.

\subsection{Cross-Scenario Validation}
\label{ssec:ctrl-sweep-validation}

The cross-scenario validation evaluates Options A and B across ten
configurations spanning the operating-point grid: dry, wet, and moderate
scenarios; 100\%, 85\%, and 70\% budgets; and $H_p = 8$ and $H_p = 14$.
Results are summarized in Table~\ref{tab:sweep-validation}.

\begin{table}[h]
\centering
\caption{Cross-scenario validation of Option~A ($\alpha_4=0,\,\alpha_6=8$)
and Option~B ($\alpha_4=2,\,\alpha_6=4$) vs.\ the nominal
$(\alpha_4=0.5,\,\alpha_6=0)$.
$\Delta y$: yield change vs.\ nominal (kg/ha);
$\Delta\mathrm{w}$: waterlog-day change (\%).}
\label{tab:sweep-validation}
\begin{tabular}{lcrr cc rr}
\toprule
\textbf{Operating point} & \textbf{Option}
  & \textbf{Yield} & $\boldsymbol{\Delta y}$
  & \textbf{wlog-d} & $\boldsymbol{\Delta\mathrm{w}}$ \\
\midrule
dry / 100\% / $H_p$=14 & A & 4{,}256 & $+8$   & 7.0  & $-68\%$ \\
dry / 100\% / $H_p$=14 & B & 4{,}256 & $+8$   & 8.7  & $-60\%$ \\
dry / 100\% / $H_p$=8  & A & 4{,}188 & $-11$  & 4.1  & $+116\%$ \\
dry / 100\% / $H_p$=8  & B & 4{,}208 & $+9$   & 5.8  & $+205\%$ \\
wet / 100\% / $H_p$=8  & A & 3{,}762 & $+161$ & 19.1 & $-47\%$ \\
wet / 100\% / $H_p$=8  & B & 3{,}707 & $+106$ & 25.4 & $-29\%$ \\
moderate / 100\% / $H_p$=14 & A & 4{,}140 & $+38$  & 3.9 & $-66\%$ \\
moderate / 100\% / $H_p$=14 & B & 4{,}143 & $+41$  & 3.6 & $-68\%$ \\
wet / 85\% / $H_p$=14 & A  & 3{,}773 & $+372$ & 24.7 & $-62\%$ \\
wet / 70\% / $H_p$=14 & A  & 3{,}768 & $+200$ & 16.0 & $-68\%$ \\
\bottomrule
\end{tabular}
\end{table}

The validation supports four conclusions. \textbf{Issue~1} (FC overshoot at
$H_p = 14$) is fully resolved: at dry/100\%/$H_p=14$, both options reduce
waterlog-days from 21.6 to 7--9 (60--68\% reduction) with yield preserved
or improved.

\textbf{Issue~3} (chronic waterlogging at wet/$H_p = 14$) is resolved
across all budgets. Option~A reduces waterlog-days from 65 to 25 at the 85\%
budget and from 50 to 16 at the 70\% budget. Notably, the wet/70\%/Option~A
yield (3{,}768~kg/ha) exceeds the nominal full-budget wet/100\% yield
(3{,}392~kg/ha) by 376~kg/ha: resolving the cost-function pathology means a
30\% budget cut can produce \emph{higher} yield than the unconstrained
original controller, because the recommended weights eliminate the
over-irrigation that was damaging the crop.

\textbf{An additional pathology at wet/$H_p=8$} is also resolved. Option~A
produces $+161$~kg/ha yield, 47\% fewer waterlog-days, zero severe
waterlogging, and 57~mm less water used: a strict Pareto improvement on every
dimension.

\textbf{The perturbation at dry/$H_p=8$ is acceptable.} Both options cause
minor waterlog-day increases (to 4.1 and 5.8, respectively). Severe
waterlogging remains at zero; yield is preserved within $\pm 20$~kg/ha; and
residual waterlog-days remain within rice's documented 1--2-day tolerance
window per event~\cite{setter1997}.

Option~A is preferred on three grounds: (i) lower peak $x_1$ across all
wet operating points; (ii) elimination of severe waterlogging in 9 of 10
validation runs versus 8 of 10 for Option~B; (iii) reduction of the cost
function from six to five active terms with no agronomic penalty.

\subsection{Ceiling Test of \texorpdfstring{$\alpha_6$}{Alpha6}}
\label{ssec:ctrl-sweep-ceiling}

To verify that $\alpha_6 = 8$ does not sit at the lower edge of an
asymptotic plateau, two further configurations were evaluated at
wet/100\%/$H_p=14$: $\alpha_6 = 10$ and $\alpha_6 = 16$ (both with
$\alpha_4 = 0$). Results are in Table~\ref{tab:sweep-ceiling}.

\begin{table}[h]
\centering
\caption{Ceiling test of $\alpha_6$ at wet/100\%/$H_p=14$. Agronomic
response saturates at $\alpha_6 = 8$; increasing beyond this value
raises worst-case solve time without proportional benefit.}
\label{tab:sweep-ceiling}
\begin{tabular}{ccccccc}
\toprule
$\boldsymbol{\alpha_6}$ & \textbf{Yield} & \textbf{Water}
  & \textbf{wlog-d} & \textbf{sev-wlog} & $x_1^{\mathrm{max}}$
  & $t_{\mathrm{solve,max}}$~(s) \\
\midrule
 8 & 3{,}774 & 321.5 & 21.1 & 0.0 & 163.3 &  921 \\
10 & 3{,}766 & 321.1 & 22.4 & 0.2 & 166.9 &  971 \\
16 & 3{,}769 & 320.1 & 18.0 & 0.2 & 163.0 & 1{,}596 \\
\bottomrule
\end{tabular}
\end{table}

Increasing $\alpha_6$ from 8 to 10 produces no improvement on any metric.
Increasing to 16 yields a 3-day waterlog-day reduction but increases
worst-case solve time by 73\% and introduces 0.2~days of severe waterlogging,
a signature of the solver becoming over-constrained on edge iterations. The
ceiling test confirms that $\alpha_6^* = 8$ sits at the asymptotic plateau
of the weight-vs-performance response and is the appropriate operating value.

% =============================================================================
\section{Recommended Operating Point}
\label{sec:ctrl-recommended}
% =============================================================================

The four phases of the sensitivity analysis converge on the recommended
operating point:
\begin{equation}
\boldsymbol{\alpha}^* = (\alpha_1, \alpha_2, \alpha_3, \alpha_4, \alpha_5,
  \alpha_6)^*
= (1.0,\; 0.016,\; 0.1,\; 0,\; 0.005,\; 8)
\label{eq:alpha-star}
\end{equation}
with prediction horizon $H_p^* = 8$~days. Each value is the result of
empirical falsification:
\begin{itemize}
  \item $\alpha_2^* = 0.016$ is anchored to the Iranian domestic-base water
    tariff and validated against all four price tiers.
  \item $\alpha_4^* = 0$ is the empirically falsified outcome of the joint
    factorial: $J_{\mathrm{ponding}}$ targets the wrong state variable.
  \item $\alpha_6^* = 8$ is identified by the factorial and confirmed by the
    ceiling test.
  \item $\alpha_3^*$ and $\alpha_5^*$ are retained at nominal values
    because the sweep showed no sensitivity worth tuning.
  \item $H_p^* = 8$ reflects equivalent yields across horizons at
    $\boldsymbol{\alpha}^*$ combined with the $5\times$ solve-time advantage
    at $H_p = 8$ relative to $H_p = 14$.
\end{itemize}

The cost function at $\boldsymbol{\alpha}^*$ resolves three pathologies
identified in the original formulation: the FC-overshoot regime at $H_p = 14$
(Issue~1), the chronic waterlogging in wet-year scenarios across all budgets
(Issue~3), and a previously unrecognized waterlogging signature at
wet/$H_p = 8$. The five-term formulation is the cost function used for the
deployed MPC throughout the remainder of the thesis and serves as the
foundation for the SAC reward design presented in the next section.

% =============================================================================
\section{Soft Actor-Critic Controller}
\label{sec:ctrl-sac}
% =============================================================================

The SAC controller addresses the principal practical limitation of the MPC:
the per-step computational burden of the online NLP. At $H_p = 8$, MPC mean
solve times are on the order of 25~s per day; at $H_p = 14$, they reach
292~s per day on average with worst-case latencies exceeding three hours in
wet scenarios. After offline training, SAC produces a fixed-parameter neural
network whose inference cost is sub-millisecond and independent of state or
forecast complexity. The training cost is amortized over many environment
interactions but pays off at every deployment step.

\subsection{Markov Decision Process Formulation}
\label{ssec:ctrl-sac-mdp}

The constrained finite-horizon optimal control problem is reformulated as a
Markov Decision Process (MDP) with the following components.

\paragraph{State $\mathbf{s}_k$ (707-dimensional).}
The state concatenates per-agent ABM observations and global scalars.
\begin{itemize}
  \item \textbf{Per-agent block} (5 $\times$ 130 = 650 entries): for each
    agent $n$, the normalized root-zone moisture $x_1^{(n)}/\mathrm{FC}$,
    normalized surface ponding $x_5^{(n)}/x_{5,\mathrm{ref}}$, normalized
    biomass $x_4^{(n)}/x_{4,\mathrm{ref}}$, raw maturation index $x_3^{(n)}$,
    and normalized elevation $\gamma^{(n)}$.
  \item \textbf{Scalar block} (9 entries, indices 650--658): day fraction
    $k/K$; remaining-budget fraction; budget-total normalization
    $W_{\mathrm{total}}/W_{\mathrm{full}}$; cumulative burn rate (water
    used per elapsed day, normalized to the daily-budget pace); today's
    rainfall; today's $\mathrm{ET}_c$; heat-stress factor $h_2$;
    cold-stress factor $h_7$; canopy growth function $g_{\mathrm{base}}$.
  \item \textbf{Forecast block} (6 $\times$ $H_p$ = 48 entries at
    $H_p = 8$, indices 659--706): $H_p$-day rolling window of rainfall,
    $\mathrm{ET}_c$, shortwave radiation, $h_2$, $h_7$, and $g_{\mathrm{base}}$.
    The forecast horizon $H_p = 8$ matches the MPC prediction horizon,
    giving both controllers identical information over the same lookahead
    window.
\end{itemize}
Total observation dimension: $650 + 9 + 48 = \mathbf{707}$.

\paragraph{Action $\mathbf{a}_k$ (130-dimensional).}
The action is a normalized irrigation vector $\mathbf{a}_k \in [0,1]^N$
with $N = 130$. The applied volume per agent is
$u^{(n)}_k = a^{(n)}_k \cdot u_{\max}$ with $u_{\max} = 12$~mm/day,
matching the MPC actuator bound. Budget compliance is enforced by the
same runner-layer clipping described in Section~\ref{sec:sys-eval}.

\paragraph{Transition function.}
The transition is the same ABM
$\mathbf{s}_{k+1} = f(\mathbf{s}_k, \mathbf{a}_k, \mathbf{w}_k;\,\mathcal{G})$
used by the MPC, ensuring that both controllers are evaluated on identical
dynamics.

\paragraph{Reward $r_k$.}
The per-step reward is constructed to approximate the negation of the MPC
path cost at $\boldsymbol{\alpha}^*$:
\begin{align}
r_k &= \;\alpha_1 \frac{\Delta\bar{x}_4(k)}{x_{4,\mathrm{ref}}}
      - \alpha_2 \frac{\sum_n u^{(n)}_k}{N \cdot u_{\max}}
      - \alpha_3 \frac{\sum_n \max(\mathrm{ST} - x_1^{(n)},\,0)}{N(\mathrm{ST}-\mathrm{WP})}
      \notag\\
    &\quad
      - \alpha_5 \frac{\sum_n (\Delta u^{(n)}_k)^2}{N \cdot u_{\max}^2}
      - \alpha_6 \frac{\sum_n \max(x_1^{(n)} - \mathrm{FC},\,0)^2}{N \cdot \mathrm{FC}^2}
      - \lambda_{\mathrm{budget}} \max(\mathrm{burn\_rate}_k - 1,\,0)^2
\label{eq:sac-reward-path}
\end{align}
where $\Delta\bar{x}_4(k) = \bar{x}_4(k) - \bar{x}_4(k-1)$ is the
field-mean daily biomass increment, and $\mathrm{burn\_rate}_k$ is the ratio
of cumulative water spent per elapsed day to the daily-budget pace. A
terminal bonus is added at episode end:
\begin{equation}
r_K^{\mathrm{terminal}} = c_{\mathrm{term}} \cdot \alpha_1 \cdot
  \frac{\bar{x}_4(K)}{x_{4,\mathrm{ref}}}
\label{eq:sac-reward-terminal}
\end{equation}
with $c_{\mathrm{term}} = 5$ ensuring that a maximum-yield episode produces
a strictly positive total return relative to the zero-action baseline. The
terminal bonus fires on \emph{both} episode termination (budget exhausted
before season end) and truncation (season ran to completion), preventing
the hoard-water local minimum that arises when termination is penalized.

The construction \eqref{eq:sac-reward-path} approximates, to within
the entropy bonus introduced by the maximum-entropy framework, the negation
of the MPC path cost at $\boldsymbol{\alpha}^*$. The approximation is exact
in the limit $\gamma \to 1$; with the operating discount $\gamma = 0.99$
over $K = 93$ steps, the effective weight on the terminal biomass term
varies between 1.0 (step 0) and $0.99^{93} \approx 0.39$ (final step),
introducing a known horizon-dependent scaling that slightly up-weights
early biomass gains relative to the MPC Mayer term. This
reward-equivalence design means that differences in MPC-vs-SAC performance
in Chapter~\ref{chap:results} are attributable to the optimization mechanism
(online NLP versus amortized neural network), not to differences in objective
function.

The budget soft-penalty coefficient $\lambda_{\mathrm{budget}} = 0.1$
provides a gradient against approaching the budget boundary. A pilot
training run with the initial value of 5.0 revealed that the penalty term
was approximately $100\times$ larger than the agronomic reward components on
any step where $\mathrm{burn\_rate} > 1$, causing entropy collapse in the
SAC optimizer (described in Section~\ref{ssec:ctrl-sac-training}). The
coefficient was calibrated to 0.1 to place the soft penalty on the same
order of magnitude as the agronomic terms while preserving the gradient
signal. The hard runner-layer budget clip guarantees constraint feasibility
independently of this soft term.

\paragraph{Discount factor.}
$\gamma = 0.99$. Because the episode is finite ($K = 93$ steps) and bounded,
the choice of $\gamma$ is less critical than in infinite-horizon problems,
but a value below 1 stabilizes the temporal-difference targets for the
critic.

\subsection{SAC Objective}
\label{ssec:ctrl-sac-objective}

The SAC algorithm~\cite{haarnoja2018sac,haarnoja2019sac} optimizes the
maximum-entropy objective:
\begin{equation}
J(\pi) = \mathbb{E}_{\pi}\!\Biggl[
  \sum_{k=0}^{K-1} \gamma^k \Bigl(
    r_k + \alpha_{\mathrm{ent}} \cdot
    \mathcal{H}\bigl[\pi(\cdot\mid\mathbf{s}_k)\bigr]
  \Bigr)
\Biggr]
\label{eq:sac-objective}
\end{equation}
where $\mathcal{H}[\pi(\cdot\mid\mathbf{s})]$ is the differential entropy
of the policy at state $\mathbf{s}$. The maximum-entropy framework produces
stochastic policies that explore effectively during training and capture
multi-modal optima — a property of value in irrigation, where several
distinct daily allocation profiles across 130 agents may achieve similar
terminal yield. The temperature $\alpha_{\mathrm{ent}}$ is tuned
automatically via a constrained dual gradient update targeting a specified
entropy level~\cite{haarnoja2019sac}.

\subsection{Gymnasium Environment Wrapper}
\label{ssec:ctrl-sac-gym}

The MDP is implemented as a Gymnasium~\cite{towers2024gymnasium}
\texttt{Env} subclass (\texttt{IrrigationEnv}) that exposes the ABM
through the standard \texttt{reset()} / \texttt{step()} interface required
by Stable-Baselines3.

\paragraph{Observation and action spaces.}
The observation space is \texttt{Box($-\infty$, $+\infty$,
shape=(707,), dtype=float32)}, matching the 707-dimensional state vector.
The action space is \texttt{Box(0, 1, shape=(130,), dtype=float32)}, a
130-dimensional unit hypercube scaled to physical irrigation volumes inside
\texttt{step()}.

\paragraph{Episode structure.}
Each episode corresponds to one 93-day Hashemi cultivation season.
\texttt{reset()} samples a training year from \texttt{TRAINING\_YEARS} and a
budget percentage uniformly from $[70\%, 100\%]$ (see
Section~\ref{ssec:ctrl-sac-training}), re-initializes the ABM to the field
transplanting conditions of Section~\ref{sec:sys-init}, and returns the
initial observation.

\texttt{step()} receives a 130-dimensional action, clips it to the
actuator bounds, scales it to irrigation volumes via $u^{(n)} = a^{(n)}
\cdot 12$~mm, applies the runner-layer budget enforcement described in
Section~\ref{sec:sys-eval}, advances the ABM by one day, and returns the
next observation, the reward \eqref{eq:sac-reward-path}, a
\texttt{terminated} flag (budget exhausted before season end), a
\texttt{truncated} flag (season ran to completion), and a diagnostic
\texttt{info} dictionary containing per-step yield, water use, and reward
decomposition.

\paragraph{Per-year precomputed quantities.}
The six biological nonlinearity arrays ($h_1, x_2, h_2, h_7, g_{\mathrm{base}},
\mathrm{ET}_c$) introduced in Section~\ref{ssec:ctrl-caching} are computed
on-the-fly for each sampled training year using that year's actual temperature
and ET data, and cached in a per-episode dictionary keyed by year integer.
This ensures that the forecast features in the SAC observation ($h_2, h_7,
g_{\mathrm{base}}, \mathrm{ET}_c$ in the scalar and forecast blocks) are
consistent with the actual climate experienced by the ABM in the same
episode. An earlier implementation used a single fixed set of precomputed
values (from the dry year, 2022) regardless of which training year was
sampled, introducing a Markov-property violation in which the observation
and the ABM transition operated on different climate data. The
per-year computation resolves this.

\subsection{CTDE Actor-Critic Architecture}
\label{ssec:ctrl-sac-ctde}

The 130-agent structure motivates a Centralized Training, Decentralized
Execution (CTDE) architecture~\cite{amato2024ctde}. CTDE separates the
actor's information set (used both during training and at deployment) from
the critic's information set (used only during training). The actor learns
a policy conditioned only on per-agent local information and a global shared
context; it is directly deployable to a distributed sensing-actuation system
without inter-agent communication at inference time. The critic accesses the
full joint state, enabling correct credit assignment across the 130 spatial
agents during training.

\paragraph{Actor: shared-parameter local policy.}
The actor is a multi-layer perceptron that takes a \textbf{62-dimensional}
input per agent:
\begin{itemize}
  \item 5 per-agent features ($x_1^{(n)}/\mathrm{FC}$, $x_5^{(n)}/x_{5,\mathrm{ref}}$,
    $x_4^{(n)}/x_{4,\mathrm{ref}}$, $x_3^{(n)}$, $\gamma^{(n)}$)
  \item 57 global features (9 scalars + 48 forecast dims)
\end{itemize}
and outputs the mean and log-standard-deviation of a Gaussian over the single-agent
action $a^{(n)}_k$. The action mean is squashed by $\tanh$ to $[0, 1]$.
The \textbf{same MLP weights are applied to all 130 agents in parallel}
through a reshape-and-broadcast operation, producing a 130-dimensional joint
action from a single forward pass. The MLP architecture is:
$62 \to 128 \to 128 \to 1$ (mean) and $62 \to 128 \to 128 \to 1$ (log-std),
with ReLU activations and log-std clamped to $[-20, 2]$.

Parameter sharing imposes the inductive bias that all agents follow the same
control policy modulo their local state and shared global context --- an
appropriate prior for a homogeneous-crop, heterogeneous-terrain field.
It reduces the actor parameter count from a naive $707 \to 130$ monolithic
MLP ($\sim$$10^6$ parameters) to a $62 \to 1$ per-agent MLP ($\sim$$10^4$
parameters), and it enforces spatial equivariance: permuting the agent index
permutes the action in the same way.

\paragraph{Critic: factorized value decomposition.}
The critic is a twin-Q value decomposition network~\cite{sunehag2018vdn}
adapted to continuous joint action spaces. Rather than mapping the full
$837$-dimensional joint state-action vector to a scalar through a monolithic
MLP, the critic decomposes the joint Q-value as a sum over per-agent local
Q-values:
\begin{equation}
Q(\mathbf{s}_k, \mathbf{a}_k)
\;=\;
\sum_{n=1}^{N}
Q_{\mathrm{loc}}\!\left(
  \mathbf{s}_k^{(n,\mathrm{loc})},\;
  \mathbf{s}_k^{(\mathrm{glob})},\;
  a_k^{(n)}
\right)
\label{eq:vdn-critic}
\end{equation}
where $\mathbf{s}_k^{(n,\mathrm{loc})}$ is the 5-dimensional local state of
agent $n$ (its own $x_1^{(n)}/\mathrm{FC}$, $x_5^{(n)}/x_{5,\mathrm{ref}}$,
$x_4^{(n)}/x_{4,\mathrm{ref}}$, $x_3^{(n)}$, $\gamma^{(n)}$),
$\mathbf{s}_k^{(\mathrm{glob})}$ is the 57-dimensional global context
(9 scalars + 48 forecast dims) broadcast identically to every agent, and
$a_k^{(n)}$ is the agent's scalar action. The local Q-function
$Q_{\mathrm{loc}}: \mathbb{R}^{63} \to \mathbb{R}$ is a $63 \to 256 \to 256
\to 1$ MLP with ReLU activations, and its parameters are shared across all
$N = 130$ agents — mirroring the parameter-sharing strategy of the actor.
Twin-Q is preserved as in standard SAC: two independent $Q_{\mathrm{loc}}$
networks $Q_{\mathrm{loc}}^{(1)}$ and $Q_{\mathrm{loc}}^{(2)}$ each produce
their own decomposed total, and the minimum is used in the Bellman
target~\cite{haarnoja2018sac}.

The factorization is motivated by the structure of the SAC reward
\eqref{eq:sac-reward-path}, in which four of the five terms
($J_{\mathrm{biomass}}$, $J_{\mathrm{water}}$, $J_{\mathrm{drought}}$,
$J_{\mathrm{overFC}}$) decompose exactly as sums over per-agent
contributions; only the budget-burn term $\lambda_{\mathrm{budget}}$
introduces global coupling, and this coupling is small because the runner-layer
budget clip enforces feasibility independently. Under such structurally
additive rewards, the linearity of the gradient with respect to the local
$Q_{\mathrm{loc}}$ contributions
($\partial Q_{\mathrm{total}}/\partial Q_{\mathrm{loc}}^{(n)} = 1$)
delivers a per-agent learning signal that the shared actor can convert into
spatially differentiated policies — restoring the spatial credit-assignment
property that the original 837-dim monolithic critic lacked. The choice and
its empirical motivation are discussed in Section~\ref{ssec:ctrl-sac-critic-motivation}.

Because the critic is used only during training, its computational cost
does not affect deployment-time inference. The total critic parameter count
is approximately 166{,}000 (two $63 \to 256 \to 256 \to 1$ MLPs), comparable
to the $\sim$540{,}000 parameters of the equivalent monolithic critic, with
the reduction arising from weight sharing across agents.

\subsection{Motivation for Value Decomposition: Spatial Credit Assignment}
\label{ssec:ctrl-sac-critic-motivation}

An initial implementation of the SAC controller used a standard monolithic
critic: the full $837$-dimensional joint state-action vector was passed
through a single $837 \to 256 \to 256 \to 1$ MLP, producing a scalar
$Q(\mathbf{s}, \mathbf{a})$. The parameter-shared actor described above was
retained unchanged. Across 200{,}000 training steps, the resulting policy
collapsed to a near-constant 2–3~mm/day irrigation rate applied uniformly to
all 130 agents, ignoring both spatial topographical variation and the
48-dimensional weather forecast. This homogenization persisted across all
five training seeds and was independent of the test scenario at evaluation
time: the policy applied 190~mm in the dry scenario and 231~mm in the wet
scenario, the opposite of the rainfall-adaptive behaviour expected of a
forecast-aware controller.

The diagnostic finding was that the monolithic critic provides no per-agent
learning signal: a single scalar Q-value cannot differentiate the relative
contribution of agent $n$'s action $a_k^{(n)}$ from that of any other agent,
so the gradient passed back to the shared actor is effectively a field-wide
average. Because all 130 agents share the same actor parameters, the only
mechanism by which the policy can produce agent-specific actions is through
the per-agent local features in the actor input — but those features only
influence the action if the critic rewards them differentially, which it does
not under the monolithic formulation. The system therefore converges to the
spatially-averaged policy that minimizes mean-field cost, regardless of how
much spatial information is present in the observation.

The decomposition \eqref{eq:vdn-critic} is the minimal architectural change
that resolves this pathology while preserving the CTDE structure: the actor
remains unchanged, the critic remains centralized at training time (each
$Q_{\mathrm{loc}}^{(n)}$ sees the full global context), and inference cost
remains sub-millisecond. The decomposition introduces no new hyperparameters
and the local MLP architecture follows the same $256 \to 256$ width as the
original monolithic critic. Empirical validation of the resulting policy
behaviour is reported in Chapter~\ref{chap:results}.

\subsection{Training Protocol}
\label{ssec:ctrl-sac-training}

\paragraph{Data split.}
A three-way split of the 26-year NASA POWER record governs the experimental
design:
\begin{itemize}
  \item \textbf{Training years (20)}: 2000, 2001, 2003--2015, 2017,
    2019--2021, 2025. Sampled uniformly at each \texttt{reset()}.
  \item \textbf{Development years (3)}: 2002 (27.1~mm, low-rainfall tercile),
    2016 (77.0~mm, mid-rainfall tercile), 2023 (88.4~mm,
    upper-mid-rainfall tercile). Used by the SB3 \texttt{EvalCallback}
    during training for \texttt{best\_model} selection and learning-curve
    monitoring. Chosen to span the training distribution by rainfall
    tercile, providing a meaningful generalization signal at each evaluation.
  \item \textbf{Test years (3)}: 2022 (dry), 2018 (moderate), 2024 (wet).
    Touched only for the final thesis comparison in
    Chapter~\ref{chap:results}. Never used during training or development-set
    evaluation.
\end{itemize}
At each training episode, the environment samples a year uniformly from the
20 training years and a budget percentage uniformly from $[70\%, 100\%]$,
exposing the policy to the full range of deficit conditions during training.

\paragraph{Hyperparameters.}
The principal training hyperparameters are summarized in
Table~\ref{tab:sac-hp}.

\begin{table}[h]
\centering
\caption{SAC training hyperparameters (final operating values).}
\label{tab:sac-hp}
\begin{tabular}{llp{6.5cm}}
\toprule
\textbf{Parameter} & \textbf{Value} & \textbf{Rationale} \\
\midrule
Total timesteps     & 500{,}000       & Pilot run converged to stable plateau by step 25{,}000; 500{,}000 provides $20\times$ safety margin \\
Replay buffer size  & 500{,}000       & Cycles $\approx$$1\times$ per run; avoids rapid stale-data churn \\
Batch size          & 256             & Standard for continuous-action SAC \\
Learning rate       & $3\!\times\!10^{-4} \to 5\!\times\!10^{-5}$ & Linear decay over total timesteps; stabilizes late-training updates \\
$\gamma$            & 0.99            & See Section~\ref{ssec:ctrl-sac-mdp} \\
$\tau$ (soft update)& 0.005           & Haarnoja et al.\ default \\
$\alpha_{\mathrm{ent}}$ & 0.05 (fixed)& Auto-tuning disabled; see entropy-coefficient note \\
$c_{\mathrm{term}}$ (terminal bonus) & 0 & Removed; see terminal-bonus note \\
$\alpha_5$ (in SAC reward)            & 0 & Removed from RL reward; see actuator-smoothing note \\
Gradient clip $\|\nabla\|_2$ & 1.0 & Caps critic update magnitude \\
Gradient steps      & 1               & $\mathrm{gs}=2$ halved throughput (37 vs.\ 68~steps/s on T4) without convergence benefit \\
Learning starts     & 1{,}000         & Random exploration before first update \\
\bottomrule
\end{tabular}
\end{table}

\paragraph{Terminal-bonus removal.}
An initial implementation used a terminal bonus
$r_K^{\mathrm{terminal}} = c_{\mathrm{term}} \alpha_1 \bar{x}_4(K)/x_{4,\mathrm{ref}}$
with $c_{\mathrm{term}} = 5$ intended to ensure that maximum-yield episodes
produced strictly positive total return. WandB telemetry from a
$200{,}000$-step pilot revealed that the resulting per-step rewards
($\mathcal{O}(10^{-2})$ to $\mathcal{O}(10^0)$) and the single terminal
spike ($\approx 5 \cdot 1 \cdot 1 = 5$) differed by approximately
$300\times$. Bellman bootstrapping of this large sparse signal back through
the 93-step horizon produced compounding overestimation in the critic,
manifested as a slow rise in $\mathcal{L}_Q$ from 0.67 (step 20{,}000) to
84{,}693 by step 130{,}000 — a 126{,}000-fold escalation that destroyed the
learned Q-function. The fix removes the terminal bonus entirely
($c_{\mathrm{term}} = 0$): the per-step biomass reward
$\alpha_1 \Delta \bar{x}_4(k)/x_{4,\mathrm{ref}}$ already integrates to the
same total biomass signal over the 93-day season when summed, distributing
the agronomic gradient signal smoothly across all timesteps rather than
concentrating it in a single backwards-propagated spike.

\paragraph{Entropy coefficient: fixed scalar.}
The standard SB3 default sets target entropy to $-\dim(\mathbf{a}) = -130$
and tunes $\alpha_{\mathrm{ent}}$ automatically via dual gradient. A pilot
run at this value produced catastrophic entropy collapse: $\alpha_{\mathrm{ent}}$
decayed from 0.17 to 0.004 within 14{,}000 steps as the agronomic reward
gradients dominated. Reducing the target entropy to $-13$ (= $-0.1 \times \dim$)
preserved early-training stability but introduced a different pathology: as
the policy converged onto a low-variance schedule in late training, the
actual policy entropy fell below the $-13$ target. The dual-gradient update
responded by spiking $\alpha_{\mathrm{ent}}$ from 0.029 (step 87{,}000) to
1.286 (step 130{,}000), injecting sudden random exploration into a critic
that had been trained to evaluate a low-variance distribution. This
distributional shock combined with the terminal-bonus bootstrap error to
produce the Phase 3 explosion of the pilot training run. The operational
fix is to bypass the dual-gradient mechanism entirely: $\alpha_{\mathrm{ent}}$
is hardcoded to a fixed scalar value of $0.05$, providing a steady,
non-adaptive exploration noise level. This eliminates the auto-tuner
distribution shock at the cost of forgoing entropy adaptation across the
training run; the loss is acceptable because the 130-dimensional action
space at hardcoded $\alpha_{\mathrm{ent}} = 0.05$ retains sufficient
exploration noise per agent throughout training.

\paragraph{Actuator smoothing in the RL reward.}
The MPC cost \eqref{eq:cost-fn} includes a $\Delta u$ regularization term
$\alpha_5 J_{\Delta u}$ that penalizes step-to-step actuator changes. This
term is appropriate for the MPC, which observes the full deterministic
8-day forecast and can smoothly schedule action transitions around future
weather events. For the SAC controller, this same penalty discourages the
behavior the policy must learn — reactive switching of irrigation in
response to rainfall events — because the policy cannot foresee that a
sudden $\Delta u$ today will reduce the loss tomorrow. Empirically, the
pilot policy learned that maintaining a constant 2~mm/day rate minimized
the $\alpha_5$ penalty regardless of rainfall, and never developed
rainfall-responsive behaviour. The $\alpha_5$ coefficient in the SAC reward
\eqref{eq:sac-reward-path} is therefore set to zero. The MPC retains
$\alpha_5 = 0.005$ unchanged. This is a deliberate asymmetry between the
two controllers' objective functions, motivated by the asymmetry in their
information access: only the MPC can plan smooth actuation around future
weather.

\paragraph{Gradient clipping and learning-rate decay.}
Two additional stabilization measures complete the v2.6 training
configuration. First, gradient clipping is applied to both actor and critic
optimizer updates with $\|\nabla\|_2 \le 1$, preventing the positive-feedback
loop in which a large critic loss produces large parameter updates that
amplify estimation errors. Second, the learning rate decays linearly from
$3\!\times\!10^{-4}$ at step~0 to $5\!\times\!10^{-5}$ at step~500{,}000,
preserving early-training exploration speed while stabilizing the
near-converged weights against late-training noise.

\paragraph{Training infrastructure and convergence.}
Training is conducted using Stable-Baselines3~\cite{raffin2021sb3} on the
Gymnasium-compliant \texttt{IrrigationEnv}, with model checkpoints saved
every 50{,}000 steps and an \texttt{EvalCallback} running 9 episodes on the
development set every 25{,}000 steps. The rotating-buffer checkpoint strategy
(overwriting the single replay-buffer file at each save interval) limits disk
footprint to $\approx$3.5~GB regardless of run length, enabling deployment
on Kaggle GPU instances with a 20~GB working-directory quota. A 50{,}000-step
pilot run (seed~0, Colab T4, 17~minutes) demonstrated clean convergence: the
policy transitioned from zero full-season completions (mean episode length
$\approx$60 days, 22\% catastrophic episodes) to 100\% full-season
completions with zero catastrophic episodes by step 25{,}000, with
development-set reward stabilizing in the range $[-2.5, -1.5]$ per episode.
The final 25{,}000 steps show no degradation, confirming that 500{,}000 steps
provides sufficient margin for stable convergence.

Training is conducted across 5 independent random seeds (0--4) to
characterize policy variance. The five seeds are run in parallel on separate
Kaggle GPU notebooks. Chapter~\ref{chap:results} reports mean and standard
deviation across seeds on the test set.

% =============================================================================
\section{Meteorological Forecast Noise Model}
\label{sec:ctrl-forecast-noise}
% =============================================================================

Both the MPC and the SAC controller receive an 8-day weather forecast as
part of their information set. In the perfect-forecast mode used for the
primary comparisons, this forecast equals the true future climate. To
evaluate robustness to realistic forecast errors, a meteorological noise
model is introduced and applied to both controllers for a separate
robustness analysis in Chapter~\ref{chap:results}.

\subsection{AR(1) Multiplicative Noise Model}
\label{ssec:ctrl-noise-model}

Operational numerical weather prediction (NWP) forecast errors are not
independent across consecutive forecast issuances: short-range ensemble
systems exhibit strong day-to-day persistence of forecast errors, with daily
precipitation forecast-error autocorrelation of $\rho \approx 0.6$ at
lead times of 1--5 days~\cite{buizza2005}. A common failure of naive noise
injection is to redraw independent errors at each timestep, causing the
controller's view of (e.g.) day $d+5$ to jump randomly between consecutive
solves on days $d$, $d+1$, $d+2$ — physically unrealistic behavior that
induces erratic MPC responses because the controller treats each forecast as
if all previous information had been discarded.

The model adopted here uses first-order autoregressive (AR(1)) multiplicative
noise, providing realistic temporal persistence:
\begin{equation}
\hat{w}(t+j \mid t) = w(t+j) \cdot \bigl(1 + \varepsilon_j(t)\bigr),
\qquad j = 1, \ldots, H_p
\label{eq:noisy-forecast}
\end{equation}
where $w(t+j)$ is the true future weather, $\hat{w}(t+j \mid t)$ is the
forecast issued at time $t$ for day $t+j$, and $\varepsilon_j(t)$ is the
lead-time-dependent error:
\begin{align}
\varepsilon_j(t) &\sim \mathcal{N}(0,\, \sigma_j^2) \quad \text{marginally},
  \quad \sigma_j = \sigma_{\mathrm{base}} \sqrt{j}
  \label{eq:noise-marginal}\\
\varepsilon_j(t+1) &= \rho \cdot \varepsilon_j(t)
  + \sqrt{1 - \rho^2} \cdot \eta_j(t), \quad
  \eta_j(t) \sim \mathcal{N}(0,\, \sigma_j^2)
  \label{eq:noise-ar1}
\end{align}
The $\sqrt{j}$ scaling in $\sigma_j$ models the degradation of forecast
skill with lead time; the AR(1) update \eqref{eq:noise-ar1} preserves the
stationary marginal variance $\sigma_j^2$ while introducing temporal
persistence $\rho$.

The operating parameters are $\sigma_{\mathrm{base}} = 0.15$ (15\% error at
1-day lead) and $\rho = 0.6$, following the Buizza et al.\ ensemble
skill analysis~\cite{buizza2005}. Noise is applied multiplicatively to
rainfall and $\mathrm{ET}_c$, with the result clipped to non-negative values
(rainfall and ET are physically non-negative). A single noise process is
shared between the two variables because their forecast errors are physically
correlated through synoptic-scale weather systems: a wet disturbance
simultaneously increases rain and decreases reference ET through increased
humidity and reduced solar radiation.

\subsection{Application to MPC and SAC}
\label{ssec:ctrl-noise-application}

For the \textbf{MPC}, the NoisyForecast model replaces the perfect forecast
passed to the NLP at each receding-horizon step. The controller solves the
NLP using corrupted rainfall and $\mathrm{ET}_c$ values; the ABM transition
still uses the true climate, so the gap between expected and actual crop
response drives deviations from the perfect-forecast MPC trajectory. This is
the standard operational mode for any model-based predictive controller
deployed without perfect future weather information.

For the \textbf{SAC controller}, the noisy forecast is injected into the
48-dimensional forecast block of the 707-dimensional observation vector at
evaluation time: the same AR(1) noise is applied to the rainfall and
$\mathrm{ET}_c$ sub-vectors before they are passed to the policy network.
The ABM transition again uses the true climate. The policy receives the same
quality of corrupted information as the MPC, enabling a fair robustness
comparison. Importantly, this evaluation is performed \emph{without
retraining}: the same policy trained on perfect forecasts from 20 historical
years is evaluated under noisy forecasts, directly testing its inherent
robustness to information degradation.

To ensure that performance differences between MPC and SAC in the
robustness analysis are attributable to policy quality rather than to
different noise realizations, both controllers use the same noise seed
($\mathrm{seed} = 42$ by default), ensuring that the AR(1) noise trajectory
is identical across all comparative runs. The Chapter~\ref{chap:results}
robustness analysis reports the performance degradation of each controller
relative to its own perfect-forecast baseline, isolating the specific
contribution of forecast noise to the performance gap.

% =============================================================================
\section{Chapter Summary}
\label{sec:ctrl-summary}
% =============================================================================

This chapter has presented the complete design and calibration of the four
controller architectures evaluated in Chapter~\ref{chap:results}.

The \textbf{MPC controller} solves a constrained finite-horizon NLP at each
of the 93 daily timesteps using a five-term cost function calibrated through
a 33-configuration, four-phase sensitivity sweep. The recommended operating
point $\boldsymbol{\alpha}^* = (1.0, 0.016, 0.1, 0, 0.005, 8)$ resolves
three pathologies of the initial formulation and anchors the water-cost weight
to the Iranian domestic-base water tariff (7{,}000~Toman/m$^3$). At
$H_p^* = 8$, mean per-day solve time is 25.9~s (worst-case 274~s) with
total season wallclock of 22--51~minutes per cell.

The \textbf{SAC controller} addresses the per-step computational burden
through amortized offline training. The CTDE architecture exploits spatial
symmetry through a parameter-shared actor ($62 \to 128 \to 128 \to 1$ per
agent) that produces 130-dimensional joint actions in a single forward pass,
paired with a twin-Q value decomposition network~\cite{sunehag2018vdn}
critic that decomposes $Q(\mathbf{s},\mathbf{a}) = \sum_{n=1}^{130}
Q_{\mathrm{loc}}(\mathbf{s}^{(n,\mathrm{loc})}, \mathbf{s}^{(\mathrm{glob})},
a^{(n)})$ using a shared $63 \to 256 \to 256 \to 1$ local MLP. The
factorization restores per-agent credit assignment that a monolithic critic
cannot provide under the reward's additively-decomposable structure, and
preserves CTDE because each local Q-value still sees the full global context
at training time. The reward is constructed as an approximate negation of
the MPC path cost at $\boldsymbol{\alpha}^*$, exact under $\gamma = 1$. Training uses a 20/3/3
train/dev/test year split from the 26-year NASA POWER record, with budget
randomized over $[70\%, 100\%]$ per episode, for 500{,}000 steps across
5 independent seeds on Kaggle T4 hardware ($\approx$2~hours per seed).
After training, SAC inference cost is sub-millisecond and constant,
irrespective of state or forecast complexity.

Both controllers are evaluated under both perfect and noisy weather forecasts
in Chapter~\ref{chap:results}. The AR(1) multiplicative noise model
($\sigma_{\mathrm{base}} = 0.15$, $\rho = 0.6$) is applied identically to
both controllers, enabling a direct comparison of the robustness of model-based
receding-horizon optimization against amortized neural-network control under
realistic forecast uncertainty. The test set (dry 2022, moderate 2018,
wet 2024) $\times$ three budget levels provides a 9-cell grid on which the
four controllers are compared in Chapter~\ref{chap:results}.
